\silentchapter{4 Landskapet er dynamisk}{landskapet er dynamisk}
\prop{4}{}{Landskapet er dynamisk.}
\prop{4.1}{a}{Seleksjonstrykka er ikkje konstante over tid\footnote{Teknologisk gjeld og endra kravspesifikasjonar syter for at trykka fluktuerer konstant.}.}
\prop{4.11}{t}{Av 3.1 og 4.1: landskapet er ein dynamisk flate.}
\prop{4.111}{t}{Haugar kan stige, søkke, flytte seg, bifurkere eller smelte saman.}
\prop{4.1111}{t}{Ein posisjon som var optimalt tilpassa under eitt sett av vilkår, kan i neste augeblink vere suboptimal\footnote{Ein arkitektur som er optimal for tusen brukarar vert suboptimal for ein million; tilpassinga er lokal i både rom og skala.}.}
\prop{4.11111}{t}{Optimal tilpassing er alltid lokal i tid.}
\prop{4.11112}{o}{Ein ny teknologi flyttar grensa for det forbodne\footnote{Nye abstraksjonslag opnar forbodne regionar av kompleksitet som før krasja grunna minne-handtering.}.}
\prop{4.111121}{o}{Ein ny marknad rekonfigurerer kva som vert belønna.}
\prop{4.111122}{o}{Eit nytt materiale utvidar mengda av moglege operasjonar.}
\prop{4.111123}{t}{Kvart av desse er eit inngrep i landskapet.}
\prop{4.11113}{o}{Den rullande medianen av høgde delt på breidde fell systematisk over fem hundreår.}
\prop{4.111131}{o}{Proporsjonen 1.36 er ikkje meir ergonomisk enn 1.88: begge er innanfor menneskeleg sitteområde.}
\prop{4.111132}{o}{Det er ein form-drift, ikkje ein funksjonsdrift.}
\prop{4.2}{o}{Endringa er diskontinuerleg.\footnote{Eldredge \& Gould (1972)\footnote{Eit paradigmeskifte i programmering manifesterer seg nøyaktig slik: stase brote av brå topologiske restruktureringar.}}}
\prop{4.21}{o}{Lange periodar med stase vert avbrotne av brå topologiske skifte.}
\prop{4.211}{o}{Det typiske forløpet er: eksplosiv radiasjon inn i ein nyopna region, fylgd av gradvis konvergens mot nye attraktorar\footnote{Nye bibliotek utløyser eksplosiv radiasjon av løysingar, før dei best eigna mønstera vert standardiserte.}.}
\prop{4.2111}{t}{Episodisk diskontinuitet er ikkje brot med den underliggande dynamikken.}
\prop{4.21111}{t}{Ho er den dynamikken si synlege signatur.}
\prop{4.21112}{o}{At endringsraten er diskontinuerleg, ikkje bimodalt fordelt, presiserer dette.}
\prop{4.21113}{o}{Det modernistiske brotet etter 1900 manifesterer seg som eit rekurrensmatrisemønster der ingen periode før 1900 liknar nokon periode etter 1900.}
\prop{4.211131}{o}{Formhistoria gjentar seg ikkje på tvers av dette brotpunktet.}
\prop{4.3}{t}{Landskapet har minne.\footnote{Arthur (1994)\footnote{Den eksisterande databasestrukturen er eit minnespor som avgrensar kva for nye tenester som rasjonelt kan implementerast.}}}
\prop{4.31}{o}{Kvar realisert form etterlet eit spor som modifiserer seleksjonstrykka for den neste.}
\prop{4.311}{t}{Forma og landskapet endrar kvarandre gjensidig.}
\prop{4.3111}{t}{All formgjeving er omformgjeving: agenten startar aldri frå ein tom posisjon\footnote{Utviklaren byrjar aldri med blanke ark; rammeverket er allereie ladet med stiavhengige seleksjonstrykk.}.}
\prop{4.31111}{t}{At agenten startar frå ein posisjon, tyder at ho arvar alt denne posisjonen impliserer.}
\prop{4.311111}{t}{Ho arvar alle naboregionane som er tilgjengelege, og alle dei som er utelukka.}
\prop{4.311112}{t}{Utgangspunktet er ikkje nøytralt.}
\prop{4.31112}{t}{Stiavhengigheit er difor eit strukturelt trekk ved formhistoria, ikkje ein anomali\footnote{Valet av eit asynkront paradigme låser systemet inn i eit spor der visse synkrone mønster vert utilgjengelege utan massiv refaktorering.}.}
\prop{4.311121}{o}{At mahogni okkuperte 16 av 16 norske stolar i perioden 1825–1849, etter å ha vore fråverande i den føregåande tilgjengelege perioden, er stiavhengig seleksjonskollaps i sitt klaraste empiriske uttrykk.}
\prop{4.4}{o}{Formrommet ekspanderer kumulativt\footnote{API-økosystem ekspanderer monotont; deprecation fjerner grensesnitt, men konseptrommet bak dei veks framleis.}.}
\prop{4.41}{o}{Nye regionar vert busette; gamle vert sjeldan heilt forlatne.}
\prop{4.411}{o}{Det kumulative konvekse hylsterveolumet veks monotont over alle tilgjengelege periodar, utan ein einaste tilbakegang.}
\prop{4.4111}{t}{Eit fast funksjonsenvelop ville krevja metning.}
\prop{4.41111}{o}{Nyhetsraten fell ikkje monotont mot metning; han hentar seg inn att når nye materiale kjem inn.}
\prop{4.41112}{t}{Det tilstøytande moglege er eit aktivt ekspanderande rom, ikkje eit tømmande lager.\footnote{Kauffman (1993)\footnote{Å openkjeldekode eit bibliotek er ikkje å tømme eit rom, men å tilgjengeleggjere eksponentielt fleire tilstøytande posisjonar for eksterne agentar.}}}
\prop{4.41113}{t}{Materialstraumar er ein hovudmekanisme for landskapsendring.}
\prop{4.411131}{o}{Kvart nytt materiale opnar nye regionar og stengjer andre.}
\prop{4.411132}{o}{Materialstraumen 1500–2025 viser klare seleksjonsbølgjer.}
\prop{4.411133}{o}{Kvar av dei er ein lokal kohortevent, ikkje gradvise overgangar mellom funksjonelt overlegne alternativ.}
\prop{4.5}{o}{Når eitt seleksjonstrykk vert dominant, kollapsar landskapet mot éin attraktor\footnote{Viss berre prosesseringsfart vektast og minnekostnad ignorerast, kollapsar systemets arkitektur mot ubrukbar rigiditet.}.}
\prop{4.51}{t}{Eit einsidig trykk som køyrer over alle andre er ikkje ei forklaring av form.}
\prop{4.511}{t}{Det er eit fråvær av forklaring.}
\prop{4.5111}{t}{Diversiteten i det realiserte formrommet er ein funksjon av mangfaldet i dei aktive trykka.}
\prop{4.51111}{t}{Ein monokultur i formrommet er eit symptom på at eitt trykk har vorte uforholdsmessig dominant.}

